\section{Conclusion}\label{sec:conclusion-new}

$\cxor$ is a particularly small member of the broader CTA family: each binary $m$-bit context
addresses one bit of state, and an XOR emitter combines that state with the source.  Under the
historically implemented $\W5$ update, the state is a hard prediction equal to the successor most
recently observed for that context.  The transform can therefore suppress stable
context-to-successor associations, producing a zero-rich residual on favourable inputs.  It
remains a length-preserving bijection, not a compressor; any coding benefit must come from a
downstream model and coder.

Varying the one-bit update exposes a complete sixteen-rule space within this architecture.
Causal sequential decoding works for every rule.  Within each interleaved context class, the
eight affine rules collapse to the operators $1$, $1+T$, and $(1+T)^{-1}$, while the eight
nonlinear rules use an observation-dependent retention factor.  The same factor gives four
adaptation horizons.  $\W5$ is the immediately adapting last-observation predictor; $\W6$ is its
running-parity operator counterpart and never forgets its initial state.  Gray coding appears
cleanly at $m=0$, but context-addressed prediction is the general $\cxor$ construction.

Moving upward from one bit to counts, probabilities, PPM, and context mixing sacrifices this tiny
classification but preserves the central principle: a causal predictor can be paired with a
reversible residual.  In this sense, $\cxor$ provides an analytically complete base case for
richer predictive models.

\paragraph{Code availability.}
The sources of this paper, the reference implementation of
Appendix~\ref{sec:implementation-new}, the historical \code{cta.py}, and the verification
harnesses of Appendix~\ref{sec:verification-new} are available in the accompanying
repository~\cite{cxorrepo} at \url{https://github.com/ewiger/cxor-paper}.
